% Irodov Part II section 2.4, original 130--139, printed pp.77--78
% (PDF pp.86--87). All entropy expressions are state functions.
\thermoentropyproblem{130}{77}
一摩尔二氧化碳气体（按理想气体处理）
绝对温度升高至原来的
$n=2.0$ 倍。
求气体熵增：
(a) 等容加热；
(b) 等压加热。
\exsolution
理想气体摩尔熵微分为
\[
\mathrm dS=C_V\frac{\mathrm dT}{T}
+R\frac{\mathrm dV}{V}
=C_p\frac{\mathrm dT}{T}-R\frac{\mathrm dp}{p}.
\]
(a) 等容 $V_2=V_1$，
\[
\boxed{\Delta S_V=C_V\ln\frac{T_2}{T_1}
=\frac{R}{\gamma-1}\ln n}.
\]
(b) 等压 $p_2=p_1$，
\[
\boxed{\Delta S_p=C_p\ln\frac{T_2}{T_1}
=\frac{\gamma R}{\gamma-1}\ln n}.
\]
按本版答案使用二氧化碳
$\gamma\simeq1.30$ 的模型值，
$n=2$，
\[
\boxed{\Delta S_V\simeq19\,\mathrm{J/K}},\qquad
\boxed{\Delta S_p\simeq25\,\mathrm{J/K}}.
\]
\excheck 二氧化碳的热容温度依赖较明显，
上述数值采用 $\gamma$ 近似常数；
更高精度应对 $C(T)/T$ 积分。

\thermoentropyproblem{131}{77}
$\nu=4.0$ mol 理想气体
在等温膨胀中熵增加
$\Delta S=23\,\mathrm{J/K}$。
求气体体积增大多少倍 $n=V_2/V_1$。
\exsolution
一摩尔理想气体等温熵变为
$R\ln(V_2/V_1)$，
故 $\nu$ 摩尔时
\[
\Delta S=\nu R\ln n.
\]
取指数：
\[
\boxed{n=\exp\left(\frac{\Delta S}{\nu R}\right)}.
\]
代入
\[
\boxed{n=\exp\left(\frac{23}{4(8.314)}\right)
\simeq2.0}.
\]
\excheck 体积比无量纲，若 $\Delta S>0$
则 $n>1$，符合等温自由膨胀或可逆膨胀的终态。

\thermoentropyproblem{132}{77}
$\nu=2.0$ mol 理想气体先等容冷却，
再等压膨胀使气体温度回到初值。
若整个过程中气压总共降低
$n=3.3$ 倍，
求系统熵变。
\exsolution
设初态 $(p_1,V_1,T_1)$。
第一段等容冷却到
$p_2=p_1/n$，
温度降至
$T_2=T_1/n$。
第二段等压升温至 $T_3=T_1$，
所以最终体积
$V_3=nV_1$，
最终压强仍为 $p_1/n$。
熵是状态函数，
由初末温度相同得
\[
\Delta S=\nu C_V\ln(T_3/T_1)
+\nu R\ln(V_3/V_1)
=\nu R\ln n.
\]
代入：
\[
\boxed{\Delta S=(2.0)R\ln(3.3)
\simeq20\,\mathrm{J/K}}.
\]
\excheck 中间过程先放热后吸热，
但熵变由初末状态确定，
不必强行把实际传热视为一个等温过程。

\thermoentropyproblem{133}{78}
氦气质量 $m=1.7\,\mathrm g$，
先进行可逆绝热膨胀，
体积增大 $n=3.0$ 倍；
随后在新压强下等压压缩，
直到气体总体积恢复初值。
求整个过程的气体熵变。
\exsolution
氦气是单原子理想气体，
$\gamma=5/3$、
$C_V=3R/2$、
$C_p=5R/2$。
设初态
$(p_1,V_1,T_1)$，
绝热膨胀后
$V_2=nV_1$，
$p_2=p_1n^{-\gamma}$。
后续等压压缩至
$V_3=V_1$，
所以末态
$p_3=p_2=p_1n^{-\gamma}$，
$T_3=T_1n^{-\gamma}$。
初末体积相同，因此
\[
\Delta S=\nu C_V\ln(T_3/T_1)
=-\nu\gamma C_V\ln n
=-\nu C_p\ln n.
\]
摩尔数
$\nu=m/M_{\rm He}=1.7/4.0=0.425$ mol，
故
\[
\boxed{\Delta S=
-\frac m{M_{\rm He}}\frac{5R}{2}\ln3
\simeq-10\,\mathrm{J/K}}.
\]
\excheck 气体系统熵可以在\emph{向环境放热}时降低；
第二定律要求整体包括环境的熵变化非负，
不要求任一子系统熵永远增加。

\thermoentropyproblem{134}{78}
$\nu=2.0$ mol 理想气体，
绝热指数 $\gamma=1.30$。
某过程中体积增大到
$\alpha=2.0$ 倍，而压强减为
$1/\beta$，$\beta=3.0$。
求气体熵变。
\exsolution
由状态方程
$T\propto pV$，
温度倍率为
$T_2/T_1=(p_2V_2)/(p_1V_1)=\alpha/\beta$。
对理想气体
\[
\Delta S=\nu C_V\ln(T_2/T_1)
+\nu R\ln(V_2/V_1).
\]
代入 $C_V=R/(\gamma-1)$：
\[
\Delta S=\frac{\nu R}{\gamma-1}
\left[\ln\left(\frac\alpha\beta\right)
+(\gamma-1)\ln\alpha\right].
\]
故
\[
\boxed{\Delta S=
\frac{\nu R}{\gamma-1}
\left(\gamma\ln\alpha-\ln\beta\right)}.
\]
给定数值：
\[
\boxed{\Delta S=
\frac{2R}{0.30}
[1.30\ln2-\ln3]\simeq-11\,\mathrm{J/K}}.
\]
\excheck 即使体积增大，若降温和降压足够显著，
系统熵变仍可能为负；
不应仅凭膨胀断定熵增加。

\thermoentropyproblem{135}{78}
容器1、2各含
$\nu=1.2$ mol 理想氦气。
两容器的容积比
$V_2/V_1=\alpha=2.0$，
绝对温度比
$T_1/T_2=\beta=1.5$。
求两种状态下氦气熵的差
$S_2-S_1$。
\exsolution
理想气体在两个平衡状态间的熵差为
\[
S_2-S_1=
\nu C_V\ln\frac{T_2}{T_1}
+\nu R\ln\frac{V_2}{V_1}.
\]
氦气为单原子气体，
$C_V=3R/2$，于是
\[
\Delta S=\nu R
\left(-\frac32\ln\beta+\ln\alpha\right).
\]
代入
$\nu=1.2,\alpha=2,\beta=1.5$：
\[
\boxed{S_2-S_1=
1.2R\left(\ln2-\frac32\ln1.5\right)
\simeq0.85\,\mathrm{J/K}\approx1.0\,\mathrm{J/K}}.
\]
\excheck 两容器各自都含相同物质的量，
可将它们看作相同气体的两个状态比较；
无需考虑不同气体混合的混合熵。

\thermoentropyproblem{136}{78}
一摩尔绝热指数为 $\gamma$ 的理想气体
经多方过程 $pV^n=\text{常数}$，
末温为初温的 $\tau$ 倍。
求该过程的熵变。
\exsolution
多方过程
$TV^{n-1}=\text{常数}$，
故
\[
\ln(V_2/V_1)=-\frac{\ln\tau}{n-1}.
\]
理想气体熵变化为
\[
\Delta S=C_V\ln(T_2/T_1)
+R\ln(V_2/V_1).
\]
代入 $C_V=R/(\gamma-1)$ 得
\[
\Delta S=\frac R{\gamma-1}\ln\tau
-\frac R{n-1}\ln\tau.
\]
整理：
\[
\boxed{\Delta S=
\frac{R(n-\gamma)}{(n-1)(\gamma-1)}
\ln\tau}.
\]
\excheck $n=\gamma$ 是可逆绝热过程，
$\Delta S=0$；
$n=1$ 是等温多方过程，此时 $\tau=1$，
若两端都是1则须单独用体积比计算熵变，
不能简单除以 $n-1$。

\thermoentropyproblem{137}{78}
$\nu=2.0$ mol 氩气在一次准静态膨胀中
压强始终与体积成正比，
$p=\alpha V$。
若气体体积增大到原来的
$r=2.0$ 倍，
求气体熵变。
\exsolution
由 $pV=\nu RT$ 和
$p\propto V$，
$T\propto V^2$，
因此
$T_2/T_1=r^2$。
单原子氩气的定容摩尔热容
$C_V=3R/2$，于是
\[
\Delta S=\nu C_V\ln(r^2)+\nu R\ln r
=\nu(2C_V+R)\ln r.
\]
因 $2C_V+R=4R$，
\[
\boxed{\Delta S=4\nu R\ln r
=8R\ln2\simeq46\,\mathrm{J/K}}.
\]
\excheck 压强与体积正比是该题的\emph{特定路径}，
不能把这一膨胀误当作等压或等温过程。

\thermoentropyproblem{138}{78}
一摩尔绝热指数为 $\gamma$ 的理想气体
沿 $p=p_0-\alpha V$
准静态变化，其中 $p_0,\alpha>0$。
求气体熵随过程变化时达到极大值所对应的体积。
\exsolution
由
$RT=pV=p_0V-\alpha V^2$，
沿过程
\[
\frac{\mathrm dT}{\mathrm dV}
=\frac{p_0-2\alpha V}{R}.
\]
熵微分为
$\mathrm dS=C_V\,\mathrm dT/T+
R\,\mathrm dV/V$，
所以
\[
\frac{\mathrm dS}{\mathrm dV}
=\frac{C_V(p_0-2\alpha V)}
{p_0V-\alpha V^2}+\frac R V.
\]
令导数为零并通分，
\[
C_V(p_0-2\alpha V)+R(p_0-\alpha V)=0.
\]
即
\[
(C_V+R)p_0-(2C_V+R)\alpha V=0.
\]
由于 $C_p=\gamma C_V$、
$2C_V+R=(\gamma+1)C_V$，
\[
\boxed{V_{\rm max}=\frac{\gamma p_0}{\alpha(1+\gamma)}}.
\]
\excheck 物理定义域是
$0<V<p_0/\alpha$；
因为 $\gamma/(\gamma+1)<1$，
极值位于允许区间内，且为熵最大值。

\thermoentropyproblem{139}{78}
一摩尔理想气体沿某过程满足
\[
S(T)=aT+C_V\ln T,
\]
其中 $a>0$、
$C_V$ 为气体定容摩尔热容。
已知体积为 $V_0$ 时温度为 $T_0$，
求温度随体积 $T(V)$ 的变化规律。
\exsolution
由理想气体的平衡态熵表达式
\[
S(T,V)=C_V\ln T+R\ln V+\text{常数},
\]
将其与题给 $S(T)=aT+C_V\ln T$ 比较，
两边的 $C_V\ln T$ 抵消后，
有
\[
R\ln V=aT+\text{常数}.
\]
用初始条件 $V=V_0,T=T_0$ 消去常数：
\[
R\ln\frac V{V_0}=a(T-T_0).
\]
因此
\[
\boxed{T(V)=T_0+\frac Ra\ln\frac V{V_0}}.
\]
\excheck 对 $a>0$，体积增大时温度上升；
模型需限制在 $T(V)>0$ 的物理定义域内。
